\documentclass[10pt]{article}
 \usepackage[margin=1in]{geometry} 
\usepackage{amsmath,amsthm,amssymb,amsfonts,color, titling}
 \usepackage{listings}

\setlength{\droptitle}{-20mm} 

\usepackage[colorlinks=true,linkcolor=red,urlcolor=blue]{hyperref}

\title{Pre-class assignment \#16}
\author{PHY-905-003\\Computational Astrophysics and
  Astrostatistics\\Spring 2026}
 \date{} % leave blank to have no date

\begin{document}
 
\maketitle

\vspace{-5mm}

\noindent
\textbf{This assignment is due the evening of Wednesday March 11, 2026.}
 Turn in all materials via GitHub Classroom.

\vspace{5mm}

\noindent
\textbf{Reading/videos:}

\begin{enumerate}

\item Read this overview of
  \href{https://www.freecodecamp.org/news/compiled-versus-interpreted-languages/}{compiled
  vs. interpreted languages}.  Note that the article points out that
Python can be run in compiled mode (if you've ever seen a
\texttt{.pyc} file, that's compiled Python code) but this is usually
reserved for imported modules.

\item Watch this video
  \href{https://www.youtube.com/watch?v=XVcRQ6T9RHo}{introducing the
    GIL} (Python's Global Interpreter Lock).


\item Read this article explaining \href{https://realpython.com/python-gil/}{why the GIL is necessary}.

\item \textbf{Optional:} A recent Python Enhancement Proposal (PEP)
  lays out a road map for removing the GIL! If you're interested, you
  can read more
  \href{https://www.infoworld.com/article/3704248/python-moves-to-remove-the-gil-and-boost-concurrency.html}{here}.  This proposal is slated to be included with Python 3.13, which is currently in the ``pre-release'' stage.
  
\end{enumerate}

\noindent
\textbf{Your assignment:}  After you read the information and watch
the video above, summarize the key points from them and put any
questions that you might have in the accompanying \texttt{ANSWERS.md}
file

\noindent
\textbf{Handing it in:}
Include your summary and questions in \texttt{ANSWERS.md} and commit it!

\end{document}
